\subsection{Procesy}
Aby wykonać obliczenia na wielu rdzeniach procesora \textbf{jednocześnie} wykorzystujemy model w którym różne 
frakcje danych są przetwarzane przez oddzielne procesy. Dzięki temu dla dużych zbiorów danych możemy 
znacznie zwiększyć wydajnośc obliczeń \\ 
W tym celu wykorzystujemy klasę \textbf{multiprocessing.Pool} z biblioteki \textbf{multiprocessing}. 
Wykorzystujemy ją do stworzenia puli procesów które potem niezależnie wykonują funkcje na różnych frakcjach 
danych.
\paragraph{Funkcje}
Procesy wykorzystujemy do:
\begin{enumerate}
    \item Obliczenia iloczynu skalarnego - Metoda $dot\_product$ wykorzystuje pulę procesów do obliczenia iloczynów par elementów dwóch wektorów, a następnie sumuje te wyniki. Zrównoleglenie tej operacji jest korzystne, gdy mamy do czynienia z bardzo długimi wektorami.
    \item Mnożenia macierzy przez wektor - $matrix\_vector\_multiply$, każdy wiersz macierzy jest mnożony przez wektor w osobnym procesie. Dzięki temu każde takie mnożenie może być przeprowadzane równolegle, co jest szczególnie efektywne dla macierzy o dużym rozmiarze.
    \item Obliczenia normy wektora - Procesy są używane do obliczenia kwadratów poszczególnych elementów wektora, a następnie sumowanie tych wartości (także w procesach) umożliwia obliczenie pierwiastka kwadratowego z ich sumy, co daje normę wektora.
    \item Działania na wektorach i macierzach - Działania takie jak dodawanie i odejmowanie wektorów, dzielenie wektora przez skalar, czy mnożenie macierzy przez skalar, są przeprowadzane w segmentach, gdzie każdy segment jest przetwarzany przez osobny proces. 
\end{enumerate}
\paragraph{Wyzwania}
\begin{itemize}
    \item Zarządzanie procesami jest kosztowne - tworzenie i zarządzanie procesami jest droższe od wątków ze względu na większy narzut systemowy.
    \item Wymiana danych między procesami - wymaga serializacji i deserializacji danych, co może wprowadzić dodatkowe opóźnienia. 
    \item Brak korzyści dla małych danych - w przypadku małych macierzy, gdzie rozmiar nie przekracza 5 tysięcy x 5 tysięcy elementów, zarządzanie procesami i koszty komunikacji międzyprocesowej mogą przewyższać korzyści wynikające z równoległego przetwarzania,
\end{itemize}

